\documentclass[11pt]{article}
\usepackage[T1]{fontenc}
\usepackage[utf8]{inputenc}
\usepackage{amsmath,amssymb,amsthm,mathtools}
\usepackage[a4paper,margin=25mm]{geometry}
\usepackage{hyperref}
\setlength{\emergencystretch}{2em}
\newtheorem{theorem}{Theorem}[section]
\newtheorem{proposition}[theorem]{Proposition}
\newtheorem{lemma}[theorem]{Lemma}
\theoremstyle{definition}
\newtheorem{definition}[theorem]{Definition}
\newtheorem{example}[theorem]{Example}
\theoremstyle{remark}
\newtheorem{remark}[theorem]{Remark}
\newcommand{\RR}{\mathbb R}
\newcommand{\NN}{\mathbb N}
\newcommand{\QQ}{\mathbb Q}
\newcommand{\HR}{{}^{*}\mathbb R}
\newcommand{\ext}[1]{{}^{*}\!#1}
\newcommand{\U}{\mathcal U}
\DeclareMathOperator{\st}{st}
\DeclareMathOperator{\Fin}{Fin}
\title{Calculus with Infinitesimals:\\Hyperreal Numbers, Transfer and Standard Parts}
\author{Revised adaptation of \emph{Calculus without Limits} by Eason Shao}
\date{Original presentations: 2024; revised: 28 September 2026}
\begin{document}
\maketitle
\paragraph{Provenance and scope.}
This document consolidates the four overlapping sources in
\texttt{calculus-without-limits}: Presentation, Semi-Finals, Finals and Poster.
The original author is Eason Shao, St Paul's School. This is an editorial and
mathematical adaptation, not an author-approved edition. It preserves the
substantive topics while replacing invalid arguments and adding explicit
foundations and proofs. Duplicate slides, overlays and informal numerical ratings
are omitted. The original MIT notice is reproduced in Appendix~\ref{app:license}.
The original visual ``microscope'' is expressed through precise relations in
Section~\ref{sec:scales}. This is a targeted revision, not a formal certification.
\paragraph{Logical framework.}
We work in ordinary classical set theory with the axiom of choice. Real
completeness, density of rationals and irrationals, elementary properties of
absolutely convergent series, and the real intermediate value theorem are
background results. The existence of a free ultrafilter is explained below.
All statements mentioning ``every positive real'' quantify over standard reals;
quantifiers over hyperreals are explicitly distinguished. The star on a function
or set denotes its natural extension, never composition with the standard-part map.
For retrieval, retain these hypotheses and the linked proof with each statement.
\tableofcontents

\section{Why an infinitesimal is not zero}
\begin{proposition}[Archimedean property of the reals]\label{prop:archimedean}
For real \(x,y>0\), some positive integer \(n\) satisfies \(nx>y\).
In particular no nonzero real number has absolute value smaller than every
positive real number.
\end{proposition}
\begin{proof}
If \(\NN\) were bounded above in \(\RR\), let \(s=\sup\NN\). Some integer
\(n>s-1\) would give \(n+1>s\), a contradiction. Thus choose \(n>y/x\).
The second assertion also follows directly: for a nonzero real \(r\), the
positive real \(|r|/2\) is smaller than \(|r|\).
\end{proof}

For nonzero \(h\), ordinary algebra gives
\[
\frac{(x+h)^2-x^2}{h}=2x+h.
\]
Taking a real limit as \(h\to0\), or taking a standard part for a nonzero
hyperreal infinitesimal \(h\), gives the derivative \(2x\). Cancelling \(h\)
uses its nonzero value. The later standard-part operation does not assert
that \(h=0\). The title ``without limits'' refers to an alternative formulation
of calculus, not an exemption from precise definitions or an assertion that
classical limit concepts have become invalid.

\section{A concrete hyperreal field}
\begin{definition}[Free ultrafilter]\label{def:ultrafilter}
A free ultrafilter \(\U\) on \(\NN=\{1,2,\ldots\}\) is a family of subsets
containing \(\NN\), not containing the empty set, closed under finite intersections
and under passing to supersets, such that exactly one of \(A\) and
\(\NN\setminus A\) belongs to \(\U\) for every \(A\subseteq\NN\), and no
finite set belongs to \(\U\). Consequently every cofinite set belongs to \(\U\).
\end{definition}
\begin{remark}[Existence assumption]\label{rem:choice}
Start with the cofinite filter and apply Zorn's lemma to the proper filters
containing it. The union of a chain is a proper filter, so a maximal one exists.
If neither \(A\) nor its complement belonged to this maximal filter, adjoining
\(A\) would still generate a proper filter: otherwise an existing filter member
would be disjoint from \(A\), forcing its complement into the filter.
This contradicts maximality. The resulting ultrafilter cannot contain a finite
set, since it already contains its cofinite complement.
Fix one such \(\U\) throughout this document.
\end{remark}

\begin{definition}[Ultrapower]\label{def:ultrapower}
For real sequences \((x_n),(y_n)\), set
\[
(x_n)\sim(y_n)\quad\Longleftrightarrow\quad\{n:x_n=y_n\}\in\U.
\]
Define \(\HR=\RR^{\NN}/\!\sim\), with elements written \([x_n]\).
Arithmetic is coordinatewise, and
\[
[x_n]<[y_n]\quad\Longleftrightarrow\quad\{n:x_n<y_n\}\in\U.
\]
Embed \(r\in\RR\) as the constant class \([r,r,\ldots]\).
\end{definition}
\begin{proposition}[Ordered-field extension]\label{prop:field}
These operations and order are well defined and make \(\HR\) an ordered
field containing an isomorphic copy of \(\RR\).
\end{proposition}
\begin{proof}
Changing representatives on a set whose complement belongs to \(\U\) does not
change any coordinatewise identity on the intersection of the relevant
\(\U\)-large sets. This proves well-definedness and the field identities.
Exactly one of \(\{x_n<y_n\},\{x_n=y_n\},\{x_n>y_n\}\) belongs to \(\U\),
which gives trichotomy; order compatibility follows coordinatewise.
If \(X=[x_n]\ne0\), the set \(A=\{n:x_n\ne0\}\) belongs to \(\U\).
Take \(y_n=1/x_n\) on \(A\), and \(y_n=0\) elsewhere. Then \(XY=1\).
Distinct constant sequences remain distinct, so the real embedding is injective.
\end{proof}

\begin{definition}[Finite, infinitesimal and infinite]\label{def:sizes}
A hyperreal \(X\) is finite if \(|X|<M\) for some positive standard real \(M\).
It is infinitesimal if \(|X|<r\) for every positive standard real \(r\).
Zero is infinitesimal; a nonzero infinitesimal is explicitly required when dividing.
A hyperreal is infinite if \(|X|>M\) for every positive standard real \(M\).
Write \(X\approx Y\) when \(X-Y\) is infinitesimal.
\end{definition}
\begin{example}[Reciprocal infinite and infinitesimal elements]\label{ex:reciprocal}
Let \(H=[n]\) and \(\varepsilon=[1/n]\). Then \(H\) is positive infinite,
\(\varepsilon\) is positive infinitesimal, and \(H\varepsilon=1\).
Indeed, each required inequality holds eventually in \(n\), hence on a
cofinite set in \(\U\). More generally, if \(\omega\) is positive infinite,
\(1/\omega\) is positive infinitesimal. An arbitrary infinitesimal and an
arbitrary infinite element need not have product one; that relation must be imposed.
\end{example}

\section{Natural extensions and the scope of transfer}
\begin{definition}[Extensions of sets and functions]\label{def:extension}
For \(A\subseteq\RR\), define
\[
\ext A=\{[x_n]:\{n:x_n\in A\}\in\U\}.
\]
For \(f:\RR\to\RR\), define \(\ext f([x_n])=[f(x_n)]\).
For \(f:D\to\RR\), the same definition on \(\ext D\) uses representatives
in \(D\) on a \(\U\)-large set; values outside that set may be chosen arbitrarily.
Finite-arity relations and operations are extended coordinatewise in the same way.
In particular, \(\ext\NN\) and \(\ext\QQ\) are defined, and \(H=[n]\in\ext\NN\).
\end{definition}
\begin{theorem}[First-order transfer in the chosen structure]\label{thm:transfer}
Consider the real ordered field together with any specified real functions,
relations and named constants. For every first-order formula \(\varphi\) of
this structure and hyperreal parameters \([a_n^1],\ldots,[a_n^k]\),
\[
\HR\models\varphi([a_n^1],\ldots,[a_n^k])
\quad\Longleftrightarrow\quad
\{n:\RR\models\varphi(a_n^1,\ldots,a_n^k)\}\in\U,
\]
where the functions and relations on the left are their natural extensions.
In particular a sentence true in the real structure remains true in its
ultrapower. This is the form of \L o\'s's theorem used here.
\end{theorem}
\begin{proof}
Atomic formulas follow from the coordinatewise definitions. Intersections and
the ultrafilter dichotomy prove the steps for conjunction and negation.
For an existential formula, a hyperreal witness gives real coordinate witnesses
on a \(\U\)-large set by induction. Conversely, if real witnesses exist on
such a set, choose one at each index there and an arbitrary real elsewhere.
Their class is a hyperreal witness by induction. Universal quantification follows
from negation and existential quantification. Structural induction proves the formula.
\end{proof}
\begin{remark}[What is not being transferred]\label{rem:scope}
This theorem is about first-order formulas in the specified structure.
It does not allow unrestricted quantification over all subsets of \(\HR\),
nor does it introduce predicates meaning ``standard'' or ``infinitesimal''.
A fixed real function, even a discontinuous one, may be named and extended.
The fact that a real function is continuous can also be expressed using
quantifiers over numbers. Quantifiers in its transferred formula range over
hyperreals, not just standard reals.
\end{remark}
\begin{example}[Arithmetic and Archimedean quantifiers]\label{ex:transfer}
Commutativity and cancellation hold in the hyperreal field. For a nonzero
infinitesimal \(\varepsilon\),
\[
\frac1{1+\varepsilon}=\frac{1-\varepsilon}{1-\varepsilon^2},
\]
since multiplying numerator and denominator by \(1-\varepsilon\) is legitimate.
The real assertion \(\forall x\in\RR\ \exists n\in\NN\ (n>x)\) transfers to
\(\forall X\in\HR\ \exists N\in\ext\NN\ (N>X)\).
It does not assert that a standard natural number exceeds every hyperreal.
Thus transfer is compatible with the infinite element \(H\).
\end{example}

\section{Standard parts and infinitesimal scales}\label{sec:scales}
\begin{lemma}[Elementary infinitesimal algebra]\label{lem:algebra}
A finite sum of infinitesimals is infinitesimal. The product of a finite
hyperreal and an infinitesimal is infinitesimal. A real infinitesimal is zero.
\end{lemma}
\begin{proof}
For a sum of \(k\ge1\) terms, bound each absolute value by \(r/k\) for any
real \(r>0\). If \(|X|<M\) with real \(M>0\), bound the other factor by
\(r/(M+1)\). For a nonzero real \(s\), use the threshold \(|s|/2\).
\end{proof}
\begin{theorem}[Existence and uniqueness of the standard part]\label{thm:st}
Every finite \(X\in\HR\) is infinitely close to exactly one real \(r\).
Define \(\st(X)=r\), with domain \(\Fin(\HR)\), the set of finite hyperreals.
\end{theorem}
\begin{proof}
The real set \(A=\{a\in\RR:a<X\}\) is nonempty and bounded above because
\(X\) is finite. Put \(r=\sup A\). For each real \(\delta>0\), some
\(a\in A\) satisfies \(r-\delta<a<X\). Also \(X<r+\delta\), since otherwise
\(r+\delta/2\in A\), contradicting the supremum. Hence \(|X-r|<\delta\).
If \(X\approx s\) for another real \(s\), then \(r-s\) is a real infinitesimal
by the triangle inequality, and is zero by Lemma~\ref{lem:algebra}.
\end{proof}
\begin{proposition}[Standard-part arithmetic]\label{prop:st-arithmetic}
For finite \(X,Y\),
\[
\st(X+Y)=\st X+\st Y,\qquad \st(XY)=(\st X)(\st Y).
\]
If \(\st Y\ne0\), then \(X/Y\) is finite and
\(\st(X/Y)=\st X/\st Y\).
\end{proposition}
\begin{proof}
Write \(X=x+u\), \(Y=y+v\), with real \(x,y\) and infinitesimal \(u,v\).
Expand the sum and product and use Lemma~\ref{lem:algebra} and uniqueness.
If \(y\ne0\), then \(|Y|>|y|/2\), so \(1/Y\) is finite, and
\(1/Y-1/y=-v/(yY)\) is infinitesimal. Apply the product formula.
\end{proof}
\begin{example}[Scales and the microscope interpretation]\label{ex:scales}
For real \(a\), nonzero infinitesimal \(\varepsilon\), and each fixed standard
integer \(k\), \(a+k\varepsilon\approx a\). Thus
\(\st(2+\varepsilon)=2\) and \(\st(3-\varepsilon+\varepsilon^4)=3\).
At an infinite scale, \(4H\), \(4H+1\) and \(4H+1+\varepsilon\) are all infinite,
so none has a real standard part. However
\((4H+1+\varepsilon)-(4H+1)=\varepsilon\), so the last two are infinitely close.
A finite drawing of these scales is schematic, not an exhaustive real coordinate
model of the hyperreal line. There is no smallest positive infinitesimal: if
\(\varepsilon>0\) is infinitesimal, so is \(\varepsilon/2\).
\end{example}

\section{Limits and continuity: all infinitesimal probes}
\begin{definition}[Classical finite limit]\label{def:limit}
Let \(D\subseteq\RR\), let \(a\) be an accumulation point of \(D\), and
let \(f:D\to\RR\). The assertion \(\lim_{x\to a}f(x)=L\in\RR\) means
\[
\forall r>0\ \exists\delta>0\ \forall x\in D:
\quad0<|x-a|<\delta\ \Longrightarrow\ |f(x)-L|<r.
\]
All quantifiers here are over real numbers.
\end{definition}
\begin{theorem}[Nonstandard criterion for a finite limit]\label{thm:limit}
In Definition~\ref{def:limit}, the limit equals \(L\) if and only if
\[
\ext f(X)\approx L\quad\hbox{for every }X\in\ext D
\hbox{ with }X\approx a\hbox{ and }X\ne a.
\]
Equivalently all these values are finite and have standard part \(L\).
\end{theorem}
\begin{proof}
Suppose the real limit holds. Fix a standard \(r>0\) and its real \(\delta>0\).
Transfer the implication with these fixed parameters. Every admissible \(X\)
is within \(\delta\) of \(a\), giving \(|\ext f(X)-L|<r\). As this holds for
all standard \(r>0\), the difference is infinitesimal.
Conversely, failure gives a real \(r_0>0\) and, for each integer \(n\), a point
\(x_n\in D\) with \(0<|x_n-a|<1/n\) and \(|f(x_n)-L|\ge r_0\).
The class \(X=[x_n]\) is infinitely close to \(a\), distinct from it, and
\(|\ext f(X)-L|\ge r_0\), contradicting the asserted criterion.
\end{proof}
\begin{proposition}[One-sided and infinite-endpoint limits]\label{prop:other-limits}
If the real domain contains a right neighbourhood of \(a\),
\(\lim_{x\downarrow a} f(x)=L\in\RR\) precisely when
\(\ext f(a+h)\approx L\) for every positive infinitesimal \(h\).
For a left neighbourhood use every negative infinitesimal \(h\).
If the domain contains a ray \((b,\infty)\), then
\(\lim_{x\to+\infty}f(x)=L\in\RR\) precisely when
\(\ext f(W)\approx L\) for every positive infinite \(W\).
For a negative ray use every negative infinite \(W\).
\end{proposition}
\begin{proof}
The one-sided proofs restrict the choices in Theorem~\ref{thm:limit} to the
appropriate side. At \(+\infty\), transfer each real tail estimate. If the
limit fails, choose \(x_n>\max\{b,n\}\) with a fixed positive error from \(L\).
Then \([x_n]\) is a positive infinite counterexample. Reflect the variable
for \(-\infty\).
\end{proof}
\begin{example}[A single infinite argument is not enough]\label{ex:single-probe}
For \(f(x)=\sin x\), the positive infinite arguments
\(W=[2\pi n]\) and \(V=[2\pi n+\pi/2]\) give respectively
\(\ext f(W)=0\) and \(\ext f(V)=1\). Thus no limit at \(+\infty\) exists,
even though the standard part at either chosen probe exists.
\end{example}
\begin{proposition}[Continuity criterion]\label{prop:continuity}
For \(a\in D\), \(f:D\to\RR\) is continuous at \(a\) if and only if
\(\ext f(X)\approx f(a)\) for every \(X\in\ext D\) with \(X\approx a\).
\end{proposition}
\begin{proof}
Transfer the epsilon--delta continuity estimate as in Theorem~\ref{thm:limit}.
If continuity fails, choose \(x_n\in D\) with \(|x_n-a|<1/n\) and a fixed
positive lower bound for \(|f(x_n)-f(a)|\). The resulting class contradicts
the criterion. At an isolated point the condition also holds: a real gap around
\(a\) transfers, forcing any such \(X\) to equal \(a\).
\end{proof}

\section{Differentiation}
\begin{theorem}[Derivative criterion]\label{thm:derivative}
Let \(f:I\to\RR\), where \(I\) is an open real interval, and let \(a\in I\).
Then \(f\) is differentiable at \(a\) with real derivative \(A\) if and only if
for every nonzero infinitesimal \(h\),
\[
\frac{\ext f(a+h)-f(a)}{h}\approx A.
\]
In this case every displayed quotient is finite and its standard part is \(A\).
\end{theorem}
\begin{proof}
Apply Theorem~\ref{thm:limit} to the real difference-quotient function
\(q(t)=(f(a+t)-f(a))/t\), defined for \(t\ne0\) with \(a+t\in I\).
Its natural extension is the displayed quotient, by coordinatewise evaluation.
Every infinitesimal \(h\) is admissible because \(a\) is an interior point.
\end{proof}
\begin{example}[Polynomial derivative]\label{ex:polynomial}
For a fixed standard integer \(m\ge1\) and \(f(x)=x^m\), the binomial formula gives
\[
\frac{(a+h)^m-a^m}{h}
=ma^{m-1}+\sum_{j=2}^m\binom mj a^{m-j}h^{j-1}.
\]
The sum is finite, with real coefficients and positive powers of an infinitesimal,
so it is infinitesimal. The criterion proves \(f'(a)=ma^{m-1}\).
For \(m=0\), the constant function has derivative zero directly.
For \(m=1\), the remainder is an empty sum. The fixed standard value of \(m\)
is essential to this finite-sum argument.
\end{example}
\begin{example}[Why the sign of the infinitesimal matters]\label{ex:abs}
At zero, \(f(x)=|x|\) gives quotient \(1\) for positive infinitesimals and \(-1\)
for negative ones. A calculation with one selected positive \(h\) therefore
cannot establish a two-sided derivative. Theorem~\ref{thm:derivative} shows
that this function has no derivative at zero.
\end{example}

\section{The exponential: retaining the required precision}
Define the real exponential by the absolutely convergent series
\(E(t)=\sum_{k=0}^{\infty}t^k/k!\), and set \(e=E(1)\).
The ratio test proves convergence for each real \(t\). The Cauchy product and
binomial theorem give \(E(x+t)=E(x)E(t)\); this uses absolute convergence.
Below \(e^x\) denotes \(E(x)\) for real \(x\); for hyperreal arguments it means
\(\ext E\), not a new unproved operation on infinitesimals.
\begin{lemma}[Quadratic error bound]\label{lem:exp-error}
For real \(|t|\le1\), \(|E(t)-1-t|\le e\,t^2\).
\end{lemma}
\begin{proof}
Absolute convergence gives
\[
|E(t)-1-t|\le\sum_{k=2}^{\infty}\frac{|t|^k}{k!}
\le t^2\sum_{k=2}^{\infty}\frac1{k!}=(e-2)t^2\le e t^2.
\]
\end{proof}
\begin{proposition}[Exponential derivative]\label{prop:exp-derivative}
For every real \(a\), \(E'(a)=E(a)\).
\end{proposition}
\begin{proof}
Transfer Lemma~\ref{lem:exp-error}. For every nonzero infinitesimal \(h\),
\[
\left|\frac{\ext E(h)-1}{h}-1\right|\le e|h|,
\]
so this quotient is infinitely close to one. The transferred addition law gives
\[
\frac{\ext E(a+h)-E(a)}h
=E(a)\frac{\ext E(h)-1}h\approx E(a).
\]
Apply Theorem~\ref{thm:derivative}. No replacement of a base by an approximately
equal base under an infinite exponent is used.
\end{proof}
\begin{proposition}[The classical sequence for \(e\)]\label{prop:e-sequence}
With \(e=\sum_{k=0}^{\infty}1/k!\), one has
\(\lim_{n\to\infty}(1+1/n)^n=e\). Consequently, for every positive infinite
hypernatural \(N\in\ext\NN\),
\[
\st\big((1+1/N)^N\big)=e.
\]
The hypernatural power denotes the natural extension of the integer-power function.
\end{proposition}
\begin{proof}
Expand
\[
(1+1/n)^n=\sum_{k=0}^n\frac1{k!}
 \prod_{j=0}^{k-1}\left(1-\frac jn\right).
\]
Every factor in each product lies in \([0,1]\), so this is at most \(e\).
For any fixed integer \(K\), the first \(K+1\) terms tend to
\(\sum_{k=0}^K1/k!\) and form a lower bound for \(n\ge K\).
Choosing \(K\) with a sufficiently small factorial-series tail proves convergence.
Transfer the resulting real tail estimates to \(N\); every infinite hypernatural
exceeds the standard threshold. The standard-part formula follows.
\end{proof}
\begin{remark}[What an approximation permits]\label{rem:precision}
The exact exponential at an infinitesimal argument is generally not \(1+h\).
For \(h>0\), transfer of \(E(t)\ge1+t+t^2/2\) for real \(t\ge0\) gives
\(\ext E(h)>1+h\). The derivative proof works because the error divided by
\(h\) is still infinitesimal, as shown above.
In general, \(X\approx Y\) alone does not permit replacing \(X\) by \(Y\)
inside a quotient by an infinitesimal: \(h\approx0\), but \(h/h=1\).
Similarly, standard-part equality of two bases is not an exact equality that
can be substituted under a variable or infinite exponent. Indeed,
\(1+1/H\approx1\), but Proposition~\ref{prop:e-sequence} gives
\(\st((1+1/H)^H)=e\), whereas \(1^H=1\). Since the series defining \(e\)
has positive terms and begins with \(1+1+1/2\), these standard parts differ.
\end{remark}

\section{The Dirichlet function is not an obstruction}
\begin{definition}[Dirichlet function and its extension]\label{def:dirichlet}
Let \(D:\RR\to\{0,1\}\) be given by \(D(x)=1\) for \(x\in\QQ\), and
\(D(x)=0\) otherwise. Its natural extension is well defined on all of \(\HR\):
\[
\ext D(X)=
\begin{cases}1,&X\in\ext\QQ,\\0,&X\notin\ext\QQ.\end{cases}
\]
\end{definition}
\begin{proof}[Verification of the extension formula]
For \(X=[x_n]\), let \(A=\{n:x_n\in\QQ\}\). Exactly one of \(A\) and its
complement lies in \(\U\). Hence \([D(x_n)]\) equals one or zero precisely
as in the displayed cases. This agrees with Definition~\ref{def:extension}.
\end{proof}
\begin{proposition}[Discontinuity is preserved]\label{prop:dirichlet}
At every real \(a\), the function \(D\) is discontinuous and has no two-sided
limit. These facts hold equally in the classical and nonstandard formulations.
\end{proposition}
\begin{proof}
By density, choose rationals \(q_n\) and irrationals \(r_n\), both distinct
from \(a\), with \(|q_n-a|<1/n\) and \(|r_n-a|<1/n\).
Then \(Q=[q_n]\) and \(R=[r_n]\) are distinct from \(a\) and infinitely close
to it, but \(\ext D(Q)=1\) and \(\ext D(R)=0\).
No real \(L\) can be infinitely close to both zero and one. Apply
Theorem~\ref{thm:limit} and Proposition~\ref{prop:continuity}.
\end{proof}
\begin{example}[Hyperrational and non-hyperrational infinitesimals]\label{ex:hyperrational}
For \(H=[n]\), the infinitesimal \(1/H=[1/n]\) belongs to \(\ext\QQ\), whereas
\(\sqrt2/H=[\sqrt2/n]\) does not, since every coordinate \(\sqrt2/n\) is
irrational. Thus \(\ext D(1/H)=1\) and \(\ext D(\sqrt2/H)=0\).
Hyperrational does not mean standard rational.
\end{example}
\begin{remark}[A different map obtained from standard parts]\label{rem:wrong-extension}
On finite hyperreals one may define \(G(X)=D(\st X)\), but this is not
\(\ext D\). For example, \(\st(\sqrt2/H)=0\), so \(G(\sqrt2/H)=D(0)=1\),
whereas \(\ext D(\sqrt2/H)=0\).
A nonstandard continuity test using this different map would not test continuity
of \(D\). The standard-part operation is not among the transferred real
operations of Theorem~\ref{thm:transfer}.
\end{remark}

\section{Intermediate values and improper integrals}
\begin{proposition}[Transferred intermediate values for a fixed function]\label{prop:ivt}
Let \(f:[a,b]\to\RR\) be a fixed continuous real function. If \(Y\in\HR\)
lies between \(f(a)\) and \(f(b)\), some \(X\in\ext{[a,b]}\) satisfies
\(\ext f(X)=Y\).
\end{proposition}
\begin{proof}
The ordinary intermediate value theorem for this fixed \(f,a,b\) is a
first-order assertion about the number \(y\) and an existential number \(x\).
Name \(f\) in the structure and apply Theorem~\ref{thm:transfer}.
This argument does not assert an intermediate value theorem for every
arbitrarily chosen map on the hyperreals.
\end{proof}
\begin{proposition}[Improper integrals through an extended primitive]\label{prop:improper}
Let \(f:[a,\infty)\to\RR\) be Riemann integrable on every compact subinterval,
and define the real function \(F(t)=\int_a^t f(x)\,dx\).
The improper integral converges to \(L\in\RR\) if and only if
\(\ext F(W)\approx L\) for every positive infinite hyperreal \(W\).
\end{proposition}
\begin{proof}
By definition convergence of this improper integral is precisely
\(\lim_{t\to+\infty}F(t)=L\). Apply Proposition~\ref{prop:other-limits}.
Thus one may use an infinite cutoff in the extended \(F\), but equality of
standard parts must hold for all such cutoffs. In particular \(W\) and \(2W\)
have the same standard part of \(\ext F\) when the integral converges.
\end{proof}
\begin{example}[A convergent and a divergent cutoff]\label{ex:integrals}
For \(f(x)=x^{-2}\) on \([1,\infty)\), the real identity \(F(t)=1-1/t\)
transfers to \(\ext F(W)=1-1/W\). Every positive infinite \(W\) gives
standard part one, so the improper integral is one.
For \(f(x)=1\), \(F(t)=t-1\) and \(\ext F(W)=W-1\) is infinite; no real
standard part exists and the improper integral does not converge to a finite value.
\end{example}

\section*{Sources and editorial status}
\addcontentsline{toc}{section}{Sources and editorial status}
The original source is Eason Shao's 2024 school presentation project
\emph{Calculus without Limits}, also titled \emph{Intuition Mathematics vs.\ Deduction
Mathematics} in its competition versions. This revision consolidates those
versions; it is not a transcription of any textbook. Definitions, hypotheses
and newly supplied proofs above are editorial work.
For a primary-author reference on the foundational framework and the universal
quantifier in the infinitesimal derivative criterion, see H. Jerome Keisler,
\emph{Foundations of Infinitesimal Calculus}, online edition dated 19 June 2022,
Chapters 1--3 and 15:
\url{https://people.math.wisc.edu/~hkeisler/foundations.pdf}.

The original bibliography consisted of Wikipedia pages on Calculus,
History of calculus, Nonstandard analysis, Hyperreal number, Infinitesimal,
Transfer principle, Limit (mathematics), Criticism of nonstandard analysis and
Influence of nonstandard analysis. These are recorded as the original author's
bibliographic context, not as independently verified authorities for this revision.
The revised mathematical statements are accompanied by proofs or explicitly
identified background assumptions. Historical date lists and subjective ratings
from the slides are not retained as mathematical corpus content.

\appendix
\section{Original MIT licence notice}\label{app:license}
\begingroup\small
MIT License

Copyright (c) 2024 Eason Shao

Permission is hereby granted, free of charge, to any person obtaining a copy
of this software and associated documentation files (the "Software"), to deal
in the Software without restriction, including without limitation the rights
to use, copy, modify, merge, publish, distribute, sublicense, and/or sell
copies of the Software, and to permit persons to whom the Software is
furnished to do so, subject to the following conditions:

The above copyright notice and this permission notice shall be included in all
copies or substantial portions of the Software.

THE SOFTWARE IS PROVIDED "AS IS", WITHOUT WARRANTY OF ANY KIND, EXPRESS OR
IMPLIED, INCLUDING BUT NOT LIMITED TO THE WARRANTIES OF MERCHANTABILITY,
FITNESS FOR A PARTICULAR PURPOSE AND NONINFRINGEMENT. IN NO EVENT SHALL THE
AUTHORS OR COPYRIGHT HOLDERS BE LIABLE FOR ANY CLAIM, DAMAGES OR OTHER
LIABILITY, WHETHER IN AN ACTION OF CONTRACT, TORT OR OTHERWISE, ARISING FROM,
OUT OF OR IN CONNECTION WITH THE SOFTWARE OR THE USE OR OTHER DEALINGS IN THE
SOFTWARE.

\endgroup
\end{document}
